% =====================================================================
% PRISM Phase 6 -- reproducibility statement + appendix (drop-in)
% Place the statement before References; place the appendix section in the
% appendix. Fills guide item 22.
% =====================================================================

% ---- Reproducibility Statement (before \bibliography) ----
\subsubsection*{Reproducibility Statement}
All benchmarks are procedurally generated, so no external data are required. The
method is fully specified in Section~\ref{sec:method} (the velocity field,
support transforms, conditional base, training objective, and the physical-space
likelihood) and the theory in Appendix~\ref{app:proofs} states every regularity
assumption. Appendix~\ref{app:repro} lists dataset sizes and splits, architecture
widths and activations, optimizer and solver settings, the Jacobian estimator,
posterior sample counts, the five random seeds, the hyperparameter-search budget,
and hardware. Every number in the paper is written to a single result manifest
and read back by the table and figure generators, so a value cannot appear twice
with two figures; the entire set of tables and figures regenerates from the
frozen manifest with one command, and all experiments re-run from scratch with
one script.

% ---- Appendix: full reproducibility details ----
\section{Reproducibility details}
\label{app:repro}

\paragraph{Data.} Each task is generated on the fly from a fixed seed. Training
uses $N_{\text{train}}$ parameter--observation pairs, validation $N_{\text{val}}$,
and test $N_{\text{test}}$, drawn i.i.d.\ from the task prior and forward process
(Table~\ref{tab:repro-data}). Because generation is procedural, simulation-based
calibration draws fresh $(\theta, y)$ pairs at evaluation time.

\paragraph{Architecture.} The liquid velocity field $g_\theta$ is a two-hidden-layer
multilayer perceptron with SiLU activations and width $W$; the amortization
network $(\boldsymbol{\mu}_\phi, \boldsymbol{\sigma}_\phi)$ and the forward head
share this width. The flow is integrated with a fixed-step Runge--Kutta method of
order $p$ using $K$ steps; the Jacobian trace uses an exact evaluation in the
low-dimensional regime and a Hutchinson estimator otherwise. Constraints use the
support transforms of Table~\ref{tab:transforms}.

\paragraph{Optimization.} All methods are trained with Adam at learning rate
$\eta$, batch size $B$, for $E$ epochs, with gradient clipping at $5.0$ and a
cosine schedule. Baselines receive method-specific tuning with an equal search
budget; the NPE-CNF baseline is validated on the analytic linear-Gaussian task
before use.

\paragraph{Evaluation.} Posterior metrics use $S$ samples per observation.
Simulation-based calibration uses $N_{\text{SBC}}$ observations. All learned-method
tables report mean $\pm$ standard deviation over seeds $\{0,1,2,3,4\}$.

\paragraph{Hardware and runtime.} Experiments run on \texttt{<GPU/CPU model>};
the full five-seed sweep completes in \texttt{<hh:mm>}. Exact commands to
regenerate each table and figure are listed in the repository \texttt{README}.

\begin{table}[h]
\centering
\caption{Dataset sizes and key hyperparameters (fill from your configuration).}
\label{tab:repro-data}
\begin{tabular}{lcccccccc}
\toprule
Task & $N_{\text{train}}$ & $N_{\text{val}}$ & $N_{\text{test}}$ & $W$ & $K$ & $\eta$ & $B$ & $E$ \\
\midrule
Darcy          & & & & & & & & \\
Burgers        & & & & & & & & \\
Helmholtz      & & & & & & & & \\
Gaussian-Linear & & & & & & & & \\
SLCP           & & & & & & & & \\
Two Moons      & & & & & & & & \\
Oscillator     & & & & & & & & \\
Lotka-Volterra & & & & & & & & \\
\bottomrule
\end{tabular}
\end{table}
